\documentclass{ximera}
%geschreven 9/11/26
%verbeterd 12/11/26
\title{Samenvatting het magnetisch veld}
\begin{document}
\begin{abstract}
\end{abstract}
\begin{onlineOnly}%voor bladindeling
    \maketitle
\end{onlineOnly}
\newgeometry{left=3cm,bottom=0.1cm}
\pagestyle{empty}

\begin{conclusion}[title={Samenvatting}]\nl

    % Permanente magneten
    \samenbox{Permanente magneten}{
        Elke magneet heeft een noordpool (N) en een zuidpool (Z). Gelijknamige polen stoten elkaar af, ongelijknamige polen trekken elkaar aan. Afzonderlijke polen bestaan niet. Door \textbf{magnetische influentie} worden \textbf{ferromagnetische} stoffen (ijzer, kobalt, nikkel) zelf magnetisch en aangetrokken. Staal is hard magnetiseerbaar, weekijzer zacht magnetiseerbaar.
    }

    % Magnetisch veld
    \samenbox{Magnetisch veld}{
        De magnetische veldvector $\vec{B}$ in een punt heeft de richting van een kompasnaald en wijst in de zin van de noordpool van de naald. De grootte wordt uitgedrukt in \textbf{tesla} (T). Veldlijnen raken aan $\vec{B}$, lopen buiten een magneet van N naar Z, vormen gesloten lussen en snijden elkaar nooit. Hoe dichter de veldlijnen, hoe sterker het veld. In een \textbf{homogeen} veld is $\vec{B}$ overal gelijk. Bij de geografische noordpool ligt een magnetische zuidpool.
    }

    % Bron
    \samenbox{Bron van het magnetisch veld}{
        Een magnetisch veld wordt veroorzaakt door \textbf{bewegende ladingen}: een stroom in een draad, of bewegende elektronen in de atomen van een magneet. Een zwaarteveld hoort bij massa's, een elektrisch veld bij ladingen en een magnetisch veld bij bewegende ladingen.
    }

    % Rechte geleider
    \samenbox{Rechte stroomvoerende geleider}{
        De veldlijnen zijn concentrische cirkels rond de draad. Zin: omklem de draad met je rechterhand, duim in de stroomzin, de vingers geven de zin van de veldlijnen. Grootte op afstand $r$, zie vergelijking~\verwijs[voor een rechte geleider]{eq:B_rechte_geleider}{magnetisme_rechte_geleider}:
        \begin{equation*}
            B = \frac{\mu_0 \cdot I}{2\pi \cdot r} \qquad \text{met } \mu_0 = 4\pi \cdot 10^{-7}\,\mathrm{T\cdot m/A}
        \end{equation*}
    }

    % Spoel
    \samenbox{Stroomvoerende spoel}{
        In de spoel is het veld homogeen, buiten de spoel lijkt het op dat van een staafmagneet. Zin: kromme vingers van de rechterhand volgens de stroom in de windingen, de duim wijst naar de noordpool. Grootte in de spoel, zie vergelijking~\verwijs[voor een spoel]{eq:B_spoel}{magnetisme_spoel}:
        \begin{equation*}
            B = \frac{\mu_0 \cdot N \cdot I}{\ell}
        \end{equation*}
    }

    % Elektromagneet
    \samenbox{Elektromagneet}{
        Een spoel met een kern van weekijzer. De kern vermenigvuldigt het veld met de relatieve permeabiliteit $\mu_r$, zie vergelijking~\verwijs[voor een spoel met kern]{eq:B_spoel_kern}{magnetisme_elektromagneet}:
        \begin{equation*}
            B = \frac{\mu_r \cdot \mu_0 \cdot N \cdot I}{\ell}
        \end{equation*}
        Een elektromagneet kan je aan- en uitschakelen, regelen en ompolen.
    }

    % Superpositie
    \samenbox{Superpositie}{
        Magnetische velden tellen vectorieel op: $\vec{B} = \vec{B}_1 + \vec{B}_2 + \dots$ Liggen de velden op één lijn, tel dan algebraïsch op. Staan ze loodrecht op elkaar, gebruik dan de stelling van Pythagoras.
    }

\end{conclusion}

\restoregeometry
\pagestyle{plain}
\end{document}
